\chapter{Resultados obtenidos}
\label{cap:resultados}

	\section{Resumen de resultados}
	\label{sec:resumen_pruebas}
	La Tabla~\ref{tab:defectos} reúne los defectos abiertos que dejaron las pruebas, con la acción
	que resuelve cada uno.



	\rowcolors{2}{rowblue}{white}
	\small
	\renewcommand{\arraystretch}{1.2}
	\setlength{\tabcolsep}{4pt}
	\begin{longtable}{|p{1.4cm}|L{5.6cm}|L{1.8cm}|L{5.7cm}|}
		\hiderowcolors\caption{Defectos abiertos y acción correctiva}
		\label{tab:defectos} \\ \showrowcolors
		\hline
		\rowcolor{headerblue}
		\headercell{Prueba} & \headercell{Defecto} & \headercell{Severidad} & \headercell{Acción} \\
		\hline
		\endfirsthead
		\hline
		\rowcolor{headerblue}
		\headercell{Prueba} & \headercell{Defecto} & \headercell{Severidad} & \headercell{Acción} \\
		\hline
		\endhead
		SEG-05 & Secretos en texto plano en el repositorio & Alta & Rotarlos, pasarlos a variables de
		entorno y depurar el historial \\ \hline
		CON-05 & Falta el encabezado HSTS & Media & Agregar \texttt{Strict-Transport-Security} en el
		servidor seguro del proxy \\ \hline
		SEG-03 & La descarga y restauración de respaldos no valida la ruta & Baja & Normalizar la ruta y
		exigir que quede dentro de la carpeta de respaldos \\ \hline
		SEG-09 & Restaurar sin confirmación responde 500 & Baja & Responder 400 desde el manejador global \\ \hline
		MAN-01 & Cuatro violaciones obligatorias de estilo & Baja & Corregir las cuatro líneas \\ \hline
	\end{longtable}
	\normalsize

\noindent Los resultados registrados respaldan el correcto funcionamiento de las
reglas cubiertas por la ejecución inicial de 364 pruebas unitarias y determinadas
integraciones verificadas en el entorno descrito. Sin embargo, las pruebas
unitarias no equivalen a una validación completa de los flujos de extremo a
extremo. La automatización con Playwright, la medición de cobertura, varias
verificaciones dinámicas de seguridad y confidencialidad, las pruebas de carga
y la aceptación formal permanecen pendientes o solo parcialmente verificadas.
Los ocho casos adicionales de resiliencia mencionados en la sección de pruebas
unitarias tampoco deben sumarse a las 364 pruebas iniciales sin incorporar su
resultado de ejecución correspondiente.

La protección de la información requiere atención prioritaria. El hallazgo
SEG-05 identifica secretos expuestos en archivos versionados; mientras no se
roten y retiren, no debe declararse completamente satisfecho el criterio de
seguridad. Asimismo, los resultados de inteligencia artificial documentan el
funcionamiento de los servicios y el comportamiento de los experimentos
sintéticos, pero no acreditan todavía calidad generalizada sobre expedientes
institucionales. La validación con expedientes reales y personal de la
Defensoría no se completó debido a las restricciones de tiempo del proyecto;
queda registrada como trabajo futuro, no como una prueba aprobada ni como una
condición que invalide las verificaciones técnicas ya realizadas. El estado
final de las pruebas debe describirse, por tanto, como una validación parcial
con resultados aprobados, hallazgos abiertos y actividades pendientes, sin
extrapolar esos resultados a condiciones que no fueron ensayadas.


\section{Evaluación experimental de los componentes de inteligencia artificial}
\label{sec:eval_ia}
La evaluación de antecedentes se organizó en consultas sintéticas con
referencias de relevancia y casos contrastivos. Se consideraron métricas de
recuperación y ordenamiento, entre ellas Precision@k, Recall@k, MRR y nDCG@k.
Se preparó una exploración de ponderaciones de RRF del 0 al 100 por ciento
para cada recuperador, manteniendo la suma de pesos constante. La configuración
final elegida para la plataforma fue 75\,\% TF-IDF y 25\,\% MiniLM.
Esta elección debe interpretarse como una decisión de compromiso: en las 200
quejas extensas alcanzó un P@5 de 52.0\,\%, próximo al 52.3\,\% de 90/10;
en las 15 consultas contrastivas alcanzó un R@10 de 13.3\,\%, por lo que
no puede afirmarse que funcione igual de bien ante paráfrasis complejas.
La selección se realizó a partir de las pruebas sintéticas disponibles y de las
necesidades de integración del prototipo, sin considerar estos resultados una
validación definitiva de su desempeño institucional. No fue posible completar
una validación con expedientes reales ni una revisión formal de los resultados
por personal de la Defensoría dentro del tiempo disponible para el Trabajo
Terminal. Asimismo, el protocolo de separación entre datos de ajuste y de
evaluación final requiere documentarse con mayor detalle para establecer
la reproducibilidad de las comparaciones. Los resultados experimentales se
presentan en el capítulo de pruebas y sus anexos.

En la verificación de infraestructura se observó un índice de 1\,210 registros.
En una actualización sin modificaciones de documentos se reutilizaron 1\,210
representaciones almacenadas y no fue necesario recalcularlas. Para el servicio
de resúmenes se ejecutó una solicitud sintética que obtuvo respuesta HTTP 200
en 26.3 segundos. Esta medición es un caso de prueba y no constituye una
garantía de latencia para cualquier entrada.

La revisión de dicha salida identificó omisiones de hechos presentes en el
texto de entrada, incluidas fechas y un cambio de calificación. Con la evidencia
recopilada no fue posible establecer si esas omisiones se produjeron durante la
selección extractiva basada en MMR, durante la reformulación con Qwen3 o por la
interacción entre ambas etapas. Una evaluación técnica por etapas puede realizarse
con narrativas sintéticas controladas comparando la entrada original, la captura
extractiva y el borrador generado. Sin embargo, la evaluación sistemática de
fidelidad y utilidad de los resúmenes para el procedimiento institucional
requeriría incorporar expedientes reales debidamente autorizados y criterios de
revisión acordados con personal de la Defensoría. Esta validación institucional no
se completó por las restricciones de tiempo del Trabajo Terminal. Por tanto,
las pruebas realizadas verifican el funcionamiento del componente e identifican
limitaciones observadas, pero no acreditan su calidad general en casos reales.


\section{Estado de avance por componente}
\label{sec:estado_avance}
La Tabla~\ref{tab:estado_avance} presenta el balance de lo implementado al cierre de esta etapa.
Se incluye de forma deliberada la columna de lo no implementado, porque el valor de este balance
para quien retome el proyecto está en lo que falta, no en lo que ya está.

\rowcolors{2}{rowblue}{white}

\small
\renewcommand{\arraystretch}{1.25}
\setlength{\tabcolsep}{4pt}

\begin{longtable}{|p{3.7cm}|p{2.1cm}|p{8.2cm}|}
	\hiderowcolors\caption{Estado de avance por componente al cierre de la etapa}
	\label{tab:estado_avance} \\ \showrowcolors
	\hline
	\rowcolor{headerblue}
	\headercell{Componente} & \headercell{Estado} & \headercell{Observaciones} \\
	\hline
	\endfirsthead

	\hline
	\rowcolor{headerblue}
	\headercell{Componente} & \headercell{Estado} & \headercell{Observaciones} \\
	\hline
	\endhead

	Portal del quejoso & Operativo & Los diecisiete requerimientos funcionales del actor están
	implementados y desplegados \\ \hline

	Panel de revisión & Operativo & Los trece requerimientos del actor de recepción están
	implementados \\ \hline

	Consola de administración & Operativo & Los nueve requerimientos de administración están
	implementados \\ \hline

	Paneles de Primer Contacto y Subdefensoría & Operativos & Disponen de interfaces institucionales; su cobertura debe contrastarse con los casos de uso definitivos \\ \hline

	Búsqueda de antecedentes & Integrada & Recuperación híbrida disponible para Primer Contacto, con revisión humana de los resultados \\ \hline

	Generación de resúmenes & Integrada, sin validación institucional & Captura extractiva y edición con Qwen3; probado con narrativas sintéticas. No se completó su evaluación con expedientes reales y personal de la Defensoría por limitaciones de tiempo; los borradores requieren revisión humana \\ \hline

	Cierre de sesión por inactividad & Pendiente & La expiración del token limita la ventana de
	exposición, pero no equivale a un cierre por inactividad \\ \hline

	Automatización integral de pruebas & Parcial & Se reportaron 364 pruebas unitarias aprobadas y se prepararon ocho casos adicionales de resiliencia; continúan pendientes la automatización integral de recorridos y verificaciones de sistema \\ \hline

	Corrección de hallazgos de análisis estático & Pendiente & Las vulnerabilidades de
	neutralización de entradas y de generación pseudoaleatoria identificadas requieren remediación
	prioritaria \\ \hline

	Separación de los cortafuegos por servidor & Pendiente & Ambos servidores comparten el mismo
	grupo de reglas; no hay exposición real hoy, pero conviene separarlos \\ \hline
\end{longtable}

\normalsize
\setlength{\tabcolsep}{6pt}
\renewcommand{\arraystretch}{1}
